\section*{Розділ VIII. Фінансування та майно}
\addcontentsline{toc}{section}{Розділ VIII. Фінансування та майно}

\subsection*{8.1. Джерела фінансування}
\addcontentsline{toc}{subsection}{8.1. Джерела фінансування}
    8.1.1. Діяльність СР СМ фінансується за рахунок коштів, передбачених у єдиному кошторисі (бюджеті) ОСС, що формується з джерел, детально визначених у Розділі V Положення про ОСС.

    8.1.2. СР СМ має право залучати благодійну допомогу, грантові та спонсорські кошти відповідно до процедур, визначених Положенням про ОСС та законодавством України.

\subsection*{8.2. Кошторис СР СМ та його виконання}
\addcontentsline{toc}{subsection}{8.2. Кошторис СР СМ та його виконання}
    8.2.1. Кошторис СР СМ є складовою частиною єдиного кошторису (бюджету) ОСС, що затверджується Конференцією студентів Університету (КСУ).

    8.2.2. Фінансовий комітет СР СМ координує збір пропозицій та розробляє проєкт кошторису СР СМ, включаючи пропозиції щодо фінансування діяльності СРГ, для подальшого включення до єдиного кошторису ОСС, що подається на затвердження КСУ.

    8.2.3. Використання коштів СР СМ здійснюється суворо відповідно до затвердженого КСУ єдиного кошторису ОСС та цільового призначення.

    8.2.4. Фінансовий комітет СР СМ здійснює загальний моніторинг виконання кошторису СР СМ. Погодження Фінансовим комітетом потребують значні або нетипові витрати, критерії яких визначаються окремим рішенням СР СМ. Поточні операційні витрати в межах затверджених статей кошторису погоджуються Головою СР СМ або відповідальним заступником Голови.

    8.2.5. Усі фінансові операції здійснюються через бухгалтерію Університету відповідно до встановлених процедур.

\subsection*{8.3. Фінансування діяльності СРГ}
\addcontentsline{toc}{subsection}{8.3. Фінансування діяльності СРГ}
    8.3.1. Фінансування діяльності СРГ здійснюється в межах коштів, передбачених для них у єдиному кошторисі ОСС, затвердженому КСУ.

    8.3.2. СРГ самостійно ініціюють та здійснюють витрати в межах затверджених для них бюджетних ліній відповідно до процедур фінансової діяльності, встановлених в Університеті.

    8.3.3. Фінансовий комітет СР СМ здійснює моніторинг відповідності витрат СРГ затвердженим бюджетним лініям та цільовому призначенню коштів, але не здійснює попереднього погодження кожної окремої транзакції в межах затверджених ліній.

\subsection*{8.4. Майно СР СМ}
\addcontentsline{toc}{subsection}{8.4. Майно СР СМ}
    8.4.1. СР СМ користується приміщеннями, обладнанням та іншим майном, наданим Адміністрацією Університету на безоплатній основі для забезпечення діяльності відповідно до Положення про ОСС.

    8.4.2. СР СМ може мати власне майно, придбане за кошти ОСС або отримане як благодійна допомога, дарунок.

    8.4.3. Матеріальну відповідальність за облік, збереження та цільове використання майна, що знаходиться у користуванні СР СМ, несе Секретар СР СМ (або інша посадова особа, визначена рішенням СР СМ відповідальною за адміністративно-господарські питання).

    8.4.4. Використання майна СР СМ здійснюється виключно для виконання статутних завдань студентського самоврядування.

\subsection*{8.5. Фінансова звітність}
\addcontentsline{toc}{subsection}{8.5. Фінансова звітність}
    8.5.1. СР СМ забезпечує прозорість використання фінансових ресурсів.

    8.5.2. Фінансовий комітет СР СМ готує та подає на затвердження СР СМ щорічний звіт про виконання кошторису СР СМ. Цей звіт включається до загального звіту СР СМ перед КСУ.

    8.5.3. Звіти про використання коштів ОСС підлягають оприлюдненню у порядку, визначеному Положенням про ОСС та рішеннями КСУ.

    8.5.4. Контроль за цільовим використанням коштів СР СМ та СРГ здійснює Студентська уповноважена делегація (СУД) відповідно до її повноважень.
